%% ============================================================
%%  Chapter 6 — Conclusions and Outlook
%%
%%  Author      : Carlos David Prado-Socorro
%%  Institution : ICMol, Universitat de València
%%
%%  Closing chapter. Synthesises the device-chemistry axis
%%  (\cref{ch:poc}, \cref{ch:bridge}, \cref{ch:comparative}) and
%%  the application axis (\cref{ch:temporal}); benchmarks against
%%  temporal / reservoir / event-driven neuromorphic hardware;
%%  states the limitations of the evidence base and the forward
%%  programme. Compiles standalone or \include'd from thesis.tex
%%  (guarded by \thesismode, as the chapters are).
%% ============================================================

%% ---- Standalone preamble (skipped when \include'd from thesis.tex) ----
\ifdefined\thesismode\else
  \documentclass[12pt, twoside]{report}
  \usepackage{thesis-format}
  \addbibresource{../bibliography/references.bib}
  \begin{document}
  \setcounter{chapter}{5}     % This is Chapter 6 when built standalone
\fi

%% --- Shared macros (\providecommand: harmless if earlier chapters declared them) ---
\providecommand{\SY}{\textit{Super Yellow}}
\providecommand{\Hybrane}{\textit{Hybrane}}
\providecommand{\PEO}{PEO}
\providecommand{\TMPE}{TMPE}
\providecommand{\LiTf}{LiOTf}
\providecommand{\ITO}{ITO}
\providecommand{\STM}{short-term memory}
\providecommand{\LTM}{long-term memory}
\providecommand{\STDP}{STDP}
\providecommand{\NARMA}{NARMA}
\providecommand{\WESAD}{WESAD}
\providecommand{\TMPE}{TMPE}
\providecommand{\twoterminal}{two-terminal}

%% ============================================================
\chapter{Conclusions and Outlook}\label{ch:conclusions}
%% ============================================================

The central question was whether the volatility of an organic polymer-electrolyte memristive device can support temporal computation. \Cref{ch:poc,ch:bridge,ch:comparative} establish the device response, diagnose its reproducibility, and measure how composition and chemistry change its dynamics. \Cref{ch:temporal} then evaluates behavioural models of those dynamics on temporal-computing tasks. The conclusions below separate the measured device evidence from the model-based results and identify the experiments needed for a hardware test.

%% ------------------------------------------------------------
\section{Summary of Contributions}\label{sec:concl_contributions}
%% ------------------------------------------------------------

The experimental and computational contributions differ in evidentiary strength because only the composition grid is replicated across devices.

\paragraph{A characterised synaptic platform (\cref{ch:poc}).} A reversible \twoterminal{} organic memristive platform combines \SY{}, \Hybrane{}, and \LiTf{} in a solution-processed active layer between \ITO{} and Ag electrodes. Across assay-specific junctions sharing that formulation and architecture, it supports analogue I--V hysteresis (an \SI{85}{\percent} conductance increase from the first to the second sweep at \SI{1}{\volt}, growing \(\sim\)\num{15}-fold over ten cycles), reversible potentiation and depression exceeding \SI{200}{\percent}, and two distinguishable retention regimes. Recalculation of the archived positive-pulse train gives \SIrange{0.084}{1.55}{\micro\joule} per event as the state changes (median \SI{1.08}{\micro\joule}). Standard-exponential fits reported in the published SI give $\tau_S\approx\SIrange{2.5}{3}{\second}$ for \STM{} and $\tau_L\approx\SI{4.7}{\second}$ for \LTM{}. The formulation also supports one numerically separated four-state excitatory post-synaptic current sequence and an asymmetric Hebbian \STDP{} window; direct re-fitting of the five-junction mean gives branch time constants of \SI{89}{\milli\second} and \SI{150}{\milli\second}. A fast-triflate/slow-lithium account is a plausible candidate interpretation, but ion motion was not measured and the later cation survey does not corroborate its simple ordering. The assays were independently peer reviewed~\cite{PradoSocorro2022}, but they do not constitute a population study; \cref{tab:ch2_assay_specimens} gives the specimen mapping.

\paragraph{A reproducibility diagnosis and a fabrication-record methodology (\cref{ch:bridge}).} Successive batches of nominally identical \Hybrane{} devices changed from low-current, high-area, multilevel behaviour to a high-conductivity response without a usable switching window. Per-device weighting, restriction to a standard protocol, sweep-amplitude matching, and recovery-tail sensitivity confirmed a multi-feature change: the on--off ratio fell from \num{2.44} to \num{1.21}, normalised hysteresis area from \num{0.235} to \num{0.027}, and potentiation at \SI{3}{\volt} changed from $+\SI{17.9}{\percent}$ to $-\SI{2.1}{\percent}$. Conduction also became more ohmic at a matched \SI{1.2}{\volt}. The change followed fabrication date rather than device age, because finished \Hybrane{} devices retained their characteristics over weeks. The evidence points to ageing of the \Hybrane{} stock, with \SY{} excluded, under a moisture-driven hydrolysis hypothesis supported by an annealing-recovery test. Replacing \Hybrane{} with poly(ethylene oxide) increased the normalised window area from \num{0.26} to \num{0.42} at \SI{3}{\volt} and maintained usable switching across roughly two and a half years. The study also produced the fabrication-record standard, normalised characterisation procedure, automated feature pipeline, and trend-analysis tools used in the later chapters.

\paragraph{Compositional and chemical control of the dynamics (\cref{ch:comparative}).} The replicated \SY{}/\PEO{}/\LiTf{}/Ag grid shows that composition tunes fading-memory time over roughly \SIrange{2}{20}{\second}; higher ion-transport-polymer content gives a weaker, faster-relaxing response. Host (\PEO{} versus \TMPE{}) and anion (triflate versus TFSI) extend the range in the illustrative chemistry survey. Cation identity also shifts the dynamics, but the triflate series runs opposite to the hypothesised \ce{Li+}>\ce{Na+}>\ce{K+} order and changes again with the anion. The available data therefore do not establish a transferable cation law. Drive amplitude changes the apparent timescale, so quantitative comparisons require a matched protocol. Equilibrium impedance falls along the composition axis and tracks fading-memory time. ATR, XRD, and UV-Vis associate the host effect with short-range ether order in an X-ray-amorphous, salt-dissolved film without a shift in the \SY{} optical gap. Because the \PEO{} fraction changes both fading memory and input nonlinearity, each composition defines a paired operating point for the simulations of \cref{ch:temporal}.

\paragraph{A temporal-computing demonstration from measured dynamics (\cref{ch:temporal}).} Compact behavioural models ($\phi \otimes \lambda \otimes f$) fitted to the \cref{ch:comparative} measurements were evaluated with linear read-outs. The corrected node update is bounded by each cell's measured potentiation peak and pulse count. Under this model, a heterogeneous composition bank changes total linear memory capacity only from \num{4.19} to \num{4.52} ($7/10$ seeds), but raises mean information-processing capacity from \num{6.31} to \num{10.33}, chiefly through nonlinear degree-two terms. Seeds quantify algorithmic sensitivity rather than experimental replication. On multi-lag \WESAD{} physiology, heterogeneity gives a modest subject-resolved gain ($R^2=0.701$ versus $0.692$ for the stronger homogeneous control). For sustained labels it does not: streaming three-class macro-F1 is $0.524$ versus $0.545$ for homogeneous memory. Under injected corruption, temporal banks exceed instantaneous input on \WESAD{} and the independent PhysioNet Non-EEG corpus, but a linear exponential-moving-average control matches or exceeds the device model. These are simulations of measured device models, not measurements from a wired reservoir.

%% ------------------------------------------------------------
\section{Answers to the Thesis Questions}\label{sec:concl_questions}
%% ------------------------------------------------------------

\paragraph{Realising and understanding an exemplar.} \emph{Can a reversible \twoterminal{} organic memristive platform made from a solution-processed polymer-electrolyte composite support analogue switching, potentiation and depression, distinct retention regimes, EPSC, and STDP responses, and can cation--anion mobility asymmetry account for the retention regimes?} The common formulation and architecture of \cref{ch:poc}, fabricated by solution processing and thermal evaporation, support these responses across assay-specific junctions. Differential displacement and relaxation of triflate and lithium are consistent with the two regimes but are not identified by the electrical measurements. The reverse and anion-dependent cation ordering in \cref{ch:comparative} means that the proposed simple mobility hierarchy is not corroborated; species-selective measurements remain necessary.

\paragraph{Composition as the quantitative control.} \emph{How do the dynamics depend on composition, and does moving the cation along the \ce{Li+}, \ce{Na+}, \ce{K+} series shift the timescale as coordination chemistry suggests?} The replicated \PEO{}/\LiTf{} grid shows a monotonic composition dependence: the ratio tunes fading-memory time over roughly an order of magnitude, with switching window and potentiation strength moving in the same direction. Cation identity shifts the dynamics, but the triflate retention series runs opposite to the predicted \ce{Li+}>\ce{Na+}>\ce{K+} order and changes again with the anion. Sample size and protocol/electrode confounds prevent a transferable cation law. Coordination chemistry remains a qualitative interpretation rather than a quantitative result.

\paragraph{Utility for temporal computing.} \emph{Can the volatile response be used for temporal computation, can composition supply task-matched timescales, and does measured variability broaden the response of a heterogeneous bank?} In the simulations of \cref{ch:temporal}, composition supplies fading-memory times that overlap peripheral affective-signal bands. A heterogeneous bank changes the lag profile and raises nonlinear information-processing capacity, but does not resolve an increase in total linear memory across seeds. The benefit therefore depends on the task and metric: the lead composition is sufficient for a single-timescale target, and the heterogeneous increment is not resolved for the sustained \WESAD{} labels.

%% ------------------------------------------------------------
\section{Benchmarking Against the Right Reference Class}\label{sec:concl_benchmark}
%% ------------------------------------------------------------

Because the device state relaxes over seconds, temporal and event-driven neuromorphic systems provide the relevant reference class~\cite{Roy2019,IndiveriLiu2015,Tanaka2019}. Non-volatile DRAM, SRAM, and crossbar ReRAM serve different storage functions. Current physical-reservoir comparisons must also separate the physical node, input encoding, output layer, and network architecture~\cite{Liang2024PhysicalRC}; task scores from different layers and protocols are not a common leaderboard.

\paragraph{Against physical reservoirs.} Physical reservoir computing has been demonstrated in dynamic memristors~\cite{Du2017,Zhong2021}, organic electrochemical devices~\cite{Cucchi2021}, and delay-coupled single nodes~\cite{Appeltant2011}. More recent systems add direct electrical tuning in two-terminal oxide devices~\cite{Lee2025TunableRC}, gate- and wavelength-controlled organic transistor dynamics~\cite{Gao2024GroupedRC}, and an organic electrochemical reservoir with a physical device read-out~\cite{Han2025ReconfigurableOECT}. These systems require fading memory and nonlinearity~\cite{Tanaka2019,Maass2002LSM}. The present devices have measured fading memory, while the fitted models have nonzero nonlinear information-processing capacity. Their sub-kilohertz response overlaps many physiological and environmental signals without time rescaling. Composition changes the fading-memory time over \SIrange{2}{20}{\second}; in the bounded simulations, the heterogeneous model bank raises nonlinear capacity but not total linear memory decisively.

\paragraph{Against organic neuromorphic devices.} Much of organic neuromorphic electronics targets linear, symmetric, non-volatile conductance updates for weight storage, as in the ENODe family~\cite{VanDeBurgt2017,VanDeBurgt2018,Fuller2019}. The devices studied here use ion-relaxation timescales for temporal processing and are compared with ionic and electrolyte-gated organic neuromorphic systems~\cite{Gkoupidenis2017,Cucchi2021}.

\begin{sidewaystable}
  \centering
  \caption[Recent reference points for tunable physical reservoirs]{Recent directly relevant reference points for tunable physical reservoirs. Temporal ranges are the values reported under each study's own extraction protocol; tasks, read-outs, and metrics differ, so the table calibrates scope rather than ranking performance. ``Physical'' means that measured device states formed the reservoir; it does not imply that every input/output operation was analogue or on-chip.}
  \label{tab:concl_sota}
  \footnotesize
  \setlength{\tabcolsep}{5pt}
  \begin{tabularx}{0.95\textheight}{@{}>{\raggedright\arraybackslash}p{3.2cm}>{\raggedright\arraybackslash}p{2.6cm}>{\raggedright\arraybackslash}p{3.3cm}>{\raggedright\arraybackslash}p{3.2cm}>{\raggedright\arraybackslash}X@{}}
    \toprule
    Study & Node and terminals & Tuning handle & Reported temporal range & Reservoir evidence and task scope \\
    \midrule
    This thesis & \SY{}/\PEO{}/\LiTf{} organic composite; two-terminal & \PEO{}/salt composition; drive amplitude & \SIrange{2}{20}{\second} across the replicated composition grid & Single-node dynamics measured; heterogeneous banks and \WESAD{}/PhysioNet tasks simulated; no wired reservoir \\
    Gao et al.\ (2024)~\cite{Gao2024GroupedRC} & Organic vertical field-effect transistor; three-terminal optoelectronic & Input wavelength and gate bias & \SIrange{0.005}{13.2}{\second} with optical input and gate bias & Physical grouped-reservoir states; image recognition and time-series prediction \\
    Han et al.\ (2025)~\cite{Han2025ReconfigurableOECT} & Solid-electrolyte OECT; three-terminal & Ethylene-glycol content selects volatile or non-volatile operation & No common fitted-$\tau$ range reported & Physical reservoir and physical device read-out; three-class human-activity classification \\
    Lee et al.\ (2025)~\cite{Lee2025TunableRC} & Pt/\ce{HfO2}/TiN memristor; two-terminal & Inter-pulse base voltage & Fitted $\tau$: \SI{0.92}{\milli\second} to \SI{123}{\milli\second}; inputs \SIrange{0.14}{100}{\milli\second} & Physical reservoir states; N-MNIST classification and Mackey--Glass prediction \\
    \bottomrule
  \end{tabularx}
\end{sidewaystable}

\paragraph{Scope of the claims.} The experimentally supported claim is that composition tunes seconds-scale fading memory in this two-terminal organic platform. The computational claim is that banks of behavioural models fitted to those measurements can support the reported benchmarks and physiological tasks. The application scores are in silico and do not establish the performance of a wired array; nor is priority claimed for a first tunable two-terminal reservoir. The simulated reservoir contains independent, non-recurrent nodes; its absolute \NARMA{} errors are high and its capacities remain below the node-count bound. Recurrent or delay-feedback connectivity could raise these values. The evidence extends to measured dynamics and model-based task validation. A fabricated, composition-diverse reservoir running the same tasks remains the missing end-to-end test.

%% ------------------------------------------------------------
\section{Limitations of the Evidence Base}\label{sec:concl_limitations}
%% ------------------------------------------------------------

The conclusions are limited by the evidence available for each part of the thesis.

\begin{enumerate}
  \item \textbf{The extended synaptic protocols are proof-of-concept rather than population-level.} EPSC, separated \STM{}/\LTM{} retention, and \STDP{} exist only for the \cref{ch:poc} \Hybrane{} formulation; some specimen identifiers are not resolved in the recovered archive, and the available multi-junction data share one substrate. They function as formulation-level lithium priors and checks in the later chapters and were not extended to comparative status across the \PEO{} family.

  \item \textbf{Only composition is replicated.} The quantitative core is the \SY{}/\PEO{}/\LiTf{}/Ag composition grid. The host, anion, and cation comparisons rest on one or two screened devices per matched cell and are illustrative. The electrode contrast is also illustrative and partly confounded with fabrication generation; its negative-polarity measurements are single-composition and uncurated, and its frequency-domain measurements rest on one batch with substrate-confounded cations. No powered material law, including a \ce{Li+}>\ce{Na+}>\ce{K+} ordering, is claimed.

  \item \textbf{Protocol-specific timescales.} Because the drive amplitude co-sets the apparent fading-memory time, absolute retention values are protocol-specific, and the supra-threshold read used in the comparative corpus differs from the sub-threshold read of the proof-of-concept device.

  \item \textbf{The composition assumption behind the models.} Pulse integration and state relaxation were characterised separately, at different write amplitudes (potentiation at \SI{3}{\volt}, decay at \SI{4}{\volt}) and one inter-pulse cadence. Combining them in the behavioural models assumes that the two regimes are compatible. The nominal composition cards already have an across-card write--retention association stronger than the central estimate transferred from the gold/\TMPE{} corpus ($\rho=+0.46$), so that estimate maps to $\kappa^{\star}=0$ in the sensitivity model. Across the tested coupling range, the heterogeneous-to-homogeneous linear-capacity ratio remains only $1.04$--$1.09$. A matched protocol is still needed for the lead composition.

  \item \textbf{The application results are in silico and only partly cross-corpus.} Every temporal-computing result is a simulation of measured device models with a linear read-out; no wired reservoir was built. The physiological-context and affective-label results use \WESAD{} as the main corpus, and heterogeneity does not improve its sustained labels. Noise robustness replicates on the PhysioNet Non-EEG corpus, but a linear moving-average control matches or exceeds the device model. Broader multi-corpus validation remains future work.

  \item \textbf{Material stability remains open.} Devices were characterised unencapsulated on a $\sim$\SI{300}{\hour} ambient shelf scale; cycling endurance, read-disturb under continuous pulsing, and long-run drift under encapsulation were not systematically studied, and device ageing was not modelled in the comparative corpus.
\end{enumerate}

%% ------------------------------------------------------------
\section{Outlook and Future Directions}\label{sec:concl_outlook}
%% ------------------------------------------------------------

The immediate experimental work concerns the model assumptions and the smaller chemistry comparisons. Stability and hardware integration require longer campaigns.

\paragraph{Matched pulse protocol.} A matched-amplitude pulse train should vary pulse count and inter-pulse interval together for one composition, electrode, and read protocol. The gold/\TMPE{} negative-polarity corpus of \cref{ch:comparative} already records integration and relaxation within single runs and finds interval-dependent potentiation and drive-depth-dependent relaxation. Repeating this measurement on the lead silver composition with a sub-threshold read would supply the missing joint parameters for the behavioural model.

\paragraph{Replicated chemistry survey.} A cation series in one host and anion, at fixed protocol and electrode, with a sub-threshold read and at least three replicates per cell, would test the \ce{Li+}>\ce{Na+}>\ce{K+} hypothesis at adequate power. The silver \TMPE{} retention follows the opposite order, \ce{K}>\ce{Na}>\ce{Li}; the gold steady-state window is flat; and the gold negative-polarity data resolve no order. The lock-in corpus contains one cation-ordered candidate: a \(\approx\!1.4\times\) \ce{K}>\ce{Na}>\ce{Li} junction-capacitance trend at one substrate per salt (\cref{ch:comparative}). Profilometry of the existing substrates and a replicated series could test whether this trend is chemical or geometric. Chemistry- and bias-resolved impedance and a generation-controlled silver-versus-gold split would strengthen the microscopic and electrode comparisons.

\paragraph{Test the degradation chemistry.} The moisture-driven hydrolytic-aging hypothesis for the \Hybrane{} stock (\cref{ch:bridge}) cannot be tested directly because the aged consumable no longer exists in its degraded state. FTIR, GPC, Karl-Fischer/TGA, and EIS on retained aged-versus-fresh stock would test the hypothesis in future material campaigns. The fabrication-record methodology now records reagent identity and history.

\paragraph{Engineer the host and salt for finer timescale control.} Modifying the polyether host and salt could extend fading memory toward the minutes-scale tonic band and separate memory time from the input nonlinearity that composition currently changes with it.

\paragraph{Stabilise for long-run inference.} Encapsulation and systematic measurements of ageing, endurance, and read disturb are required to assess continuous operation and remove uncontrolled ageing from later comparisons.

\paragraph{Build the hardware reservoir.} A wired bank of composition-diverse \PEO{}/\LiTf{} devices with an on-chip linear read-out would test the simulations directly. Integration toward a 1T1M-style array and operation on an affective stream would measure whether the predicted benefit of timescale heterogeneity survives in hardware.

\paragraph{Physiological sensor interfaces.} Solution processing, mechanical flexibility, biocompatible chemistry, and seconds-scale dynamics make these devices candidates for skin-conformal sensing and inference. In the frequency-domain measurements of \cref{ch:comparative}, a DC-biased junction acts as a voltage-programmable low-pass amplifier below \SI{10}{\hertz}. A hardware experiment could test whether one junction can combine pre-amplification and temporal filtering at a physiological sensor interface.

%% ------------------------------------------------------------
\section{Closing Remarks}\label{sec:concl_closing}
%% ------------------------------------------------------------

The replicated experiments identify composition as a control of seconds-scale fading memory in these polymer-electrolyte composites; the cation trend expected from coordination chemistry is not transferable across the available data. Simulations fitted to the measured dynamics show how mixed-composition banks can process temporal signals, but they stop short of a wired reservoir. A matched-amplitude pulse protocol, a replicated chemistry series, and a composition-diverse hardware bank would test the remaining assumptions.

%% ---- Standalone close (skipped when \include'd from thesis.tex) ----
\ifdefined\thesismode
  \let\thesisstandaloneclose\relax
\else
  \printbibliography[heading=bibintoc, title={References}]
  \def\thesisstandaloneclose{\end{document}}
\fi
\thesisstandaloneclose
